\centerline{\bf \Huge Первая МатДрака}

\textit{ВСЕ участники команды, занявшей I место, получают по 4 балла в рейтинговую таблицу и право на 4 попытки ввода кода.}

\textit{ВСЕ участники команды, занявшей II место, получают по 3 балла в рейтинговую таблицу и право на 3 попытки ввода кода.}

\textit{ВСЕ участники команды, занявшей III место, получают по 2 балла в рейтинговую таблицу и право на 2 попытки ввода кода.}

\textit{ВСЕ участники команды, занявшей IV место, получают по 1 баллу в рейтинговую таблицу и право на 1 попытку ввода кода.}

\textbf{Ученик, решивший задачу, получает 2 балла в рейтинговую таблицу и право на 2 попытки ввода кода; если решали двое, то по 1 баллу и право на 1 попытку ввода кода.}

\problem{\textbf{1}} Чему равно следующее выражение:  $2^{14}-2^{11}-2^{9}+2^{6}-979$\\

\problem{\textbf{2}} Ягами выписывает последовательно на доску по возрастанию все числа, в которых число четных цифр равно числу нечетных цифр. Какое число выпишет Ягами 46–м?\\

\problem{\textbf{3}} Сколько существует трёхзначных чисел, у которых сумма цифр больше произведения цифр?\\


\problem{\textbf{4}} Сколько на шахматной доске имеется всевозможных прямоугольников, состоящих из четырёх клеток?\\

\problem{\textbf{5}} Микаса заменила в примере на умножение двузначных чисел цифры буквами: одинаковые – одинаковыми, разные – разными. У неё получилось АБ$\times$ВГ = БББ Каким мог быть исходный пример? Найдите все возможные варианты.\\


\problem{\textbf{6}} Найдите 2009–й член последовательности: 1, 2, 1, 2, 3, 2, 1, 2, 3, 4, 3, 2, 1, 2, 3, 4, 5, 4, 3, 2, 1…



\problem{\textbf{7}} Замените звёздочки цифрами так, чтобы равенство стало верным
и все семь цифр были различными: *** + ** = 1056.\\




\problem{\textbf{8}} В саду у Леонардо и Микеланджело росло 2026 розовых кустов. Микеланджело
полил половину всех кустов, и Леонардо полил половину всех кустов. При этом оказалось, что ровно
19 кустов, самых красивых, были политы и Леонардо, и Микеланджело. Сколько розовых кустов остались
не политыми?\\


\problem{\textbf{9}} Решите ребус: БАО $\times$ БА $\times$ Б = 2002.\\


\problem{\textbf{10}} Среди 35 учеников 5 «А» класса не любят есть конфеты 13, торты —
12, пряники — 9 учеников. Кроме того, не любят есть конфеты и торты 3, пряники и торты — 6,
конфеты и пряники — 5 учеников. Наконец, не любят есть конфеты, торты и пряники 2 ученика.
Сколько в классе ребят, которые любят есть конфеты, торты и пряники?\\


\problem{\textbf{11}} В группе 17 человек знают английский язык, 14 человек знают
китайский язык, 19 человек знают корейский язык и 20 человек знают калмыцкий язык. При этом
34 человека в группе знают ровно один язык из перечисленных, а остальные — ровно два языка
из перечисленных. Сколько человек в группе?\\


\problem{\textbf{12}} Решите ребус: ABCDEF $\times$ 3 = BCDEFA.\\


\problem{\textbf{13}} Решите ребус: MSU + MSU + MSU + MSU + OLYMP + OLYMP = MOSCOW.\\


\problem{\textbf{14}} Из натурального числа вычли сумму его цифр, из
полученного числа снова вычли сумму его (полученного числа) цифр и т.д. После одиннадцати
таких вычитаний получился нуль. С какого числа начинали? Перечислите все варианты!\\


\problem{\textbf{15}} Решите ребус: 250 $\times$ ЛЕТ + МГУ = 2005 $\times$ ГОД.



